\section*{Глава 1. Анализ предметной области}
\addcontentsline{toc}{section}{Глава 1. Анализ предметной области}
\stepcounter{section}

\subsection{Анализ существующих инструментов и алгоритмов организации учебных очередей}

В современной образовательной среде эффективное распределение временных ресурсов преподавателя и 
студента является критически важным фактором. Проблема «живых» очередей на кафедрах технических вузов, 
таких как Московский Политех, остается актуальной из-за высокой плотности учебного графика и сложности принимаемых работ 
(лабораторных работ, курсовых проектов). В данном подразделе проводится подробный анализ теоретических моделей (алгоритмов)
 и практических инструментов, применяемых для решения этой задачи.

\subsubsection{Анализ существующих алгоритмов организации учебных очередей и обоснование эффективности разработки}

Организация очередности — это классическая задача теории массового обслуживания (ТМО). 
Для учебного процесса выбор алгоритма напрямую влияет на пропускную способность «канала обслуживания» 
(преподавателя) и удовлетворенность «потребителей» (студентов).

Рассмотрим основные алгоритмические подходы:

\begin{itemize}
 \item \textbf{Алгоритм FIFO (First-In-First-Out):}
 \begin{itemize}
  \item \textit{Механика:} Студенты добавляются в конец списка по мере готовности или прихода в аудиторию.
  \item \textit{Преимущества:} Простота реализации и интуитивная справедливость.
  \item \textit{Критический недостаток:} Возникновение «блокирующих элементов». Если первый студент в 
  очереди защищает сложный проект в течение 40 минут, остальные участники теряют мотивацию и время. 
  В условиях пары (90 минут) это приводит к тому, что за одно занятие успевают защититься всего 2–3 человека.
 \end{itemize}
 
 \item \textbf{Алгоритм SJF (Shortest Job First):}
 \begin{itemize}
  \item \textit{Применение в учебе:} Преподаватель сначала принимает короткие отчеты или исправленные замечания, 
  а затем переходит к комплексным защитам.
  \item \textit{Эффективность:} Максимально быстрое уменьшение длины очереди.
  \item \textit{Недостаток:} Студенты со сложными работами могут бесконечно отодвигаться в конец списка («эффект голодания»).
 \end{itemize}
 
 \item \textbf{Круговое обслуживание (Round Robin):}
 \begin{itemize}
  \item \textit{Механика:} Каждому студенту выделяется фиксированный лимит времени (тайм-слот), например, 
  10–15 минут. Если работа не защищена за отведенное время, студент перемещается в конец очереди, 
  а преподаватель переходит к следующему.
  \item \textit{Обоснование:} Данный алгоритм является наиболее социально справедливым в больших группах, 
  так как гарантирует внимание каждому студенту в течение одной пары.
 \end{itemize}
\end{itemize}

\textbf{Обоснование разработки:} Проектируемое приложение позволит реализовать 
гибридный адаптивный алгоритм. Система будет использовать LIFO как базовый метод.
Использования метода LIFO в кооперации с дополнительными возможностями описанными позже, подразумевает то, что
 запись происходит последовательно по времени, кто раньше записался, тот поднимается вверх по списку в очереди,
 далее с этого списка уходит в самых конец если захочет сдавать еще раз. Так же о доп возможностях, приложение
  предоставит преподавателю и администратору(в данной работе это роль отдается старосте группы) 
инструменты динамического управления (возможность передвигать, заменять или убирать студентов из очереди, 
так же быстрая смена студентов во время самого занятия). Это позволит минимизировать простои и исключить 
конфликты, связанные с субъективным восприятием очередности.

\subsubsection{Анализ существующих инструментов организации учебных очередей и обоснование эффективности разработки}

\begin{table}[ht!]
    \centering
    \caption{Сравнительный анализ инструментов организации учебной очереди}
    \label{tab:queue_tools}
    \small % Немного уменьшим шрифт, чтобы влезло больше данных
    \begin{tabularx}{\textwidth}{l|p{2cm}|p{2cm}|p{2cm}|p{2.5cm}|p{2.5cm}}
        \toprule
        \textbf{Критерий} & \textbf{Google Таблицы} & \textbf{Telegram-боты} & \textbf{CRM-системы (YCLIENTS и др.)} & \textbf{Предлагаемое решение} \\
        \midrule
        \textbf{Учёт успеваемости} & Нет & Нет & Нет & \textbf{Да} (кол-во сданных работ) \\
        \midrule
        \textbf{Гибкое время} & Нет (только текст) & Нет & Только жесткие слоты & \textbf{Да} (раньше/позже/все равно) \\
        \midrule
        \textbf{Прогноз длительности} & Нет & Нет & Фикс. время услуги & \textbf{Да} (на основе типа работы) \\
        \midrule
        \textbf{Подтверждение записи} & Нет & Редко & SMS/Push (платно) & \textbf{Да} (за 24ч до пары) \\
        \midrule
        \textbf{Защита от правок} & Отсутствует & Средняя & Высокая & \textbf{Высокая} (ролевая модель) \\
        \midrule
        \textbf{Специализация} & Универсально & Общая & Коммерция & \textbf{Академическая} \\
        \bottomrule
    \end{tabularx}
\end{table}

\textbf{Детальный разбор аналогов:}

\begin{itemize}
	\item \textbf{Физическая очередь:} Основная проблема — «эффект присутствия». 
	Студент вынужден находиться в коридоре вуза в течение 4–6 часов, что не позволяет эффективно 
	использовать время для самоподготовки в библиотеке или коворкинге.
	
	\item \textbf{Общие облачные документы (Google Sheets):} Являются «стандартом де-факто» во многих группах
	 Московского Политеха. Однако отсутствие жесткой логики приложения позволяет любому пользователю стереть 
	 чужую фамилию или изменить порядок записей. Это создает почву для академических конфликтов. Кроме того, 
	 интерфейс таблиц плохо адаптирован для мобильных устройств в условиях нестабильного Wi-Fi.
	
	\item \textbf{Telegram-боты:} Хорошо справляются с задачей регистрации, но обладают линейным интерфейсом. 
	Преподавателю крайне неудобно пролистывать историю чата, чтобы увидеть общую картину загруженности на день. 
	Отсутствует возможность визуализации структуры базы данных (например, какие лабораторные работы чаще всего 
	вызывают затруднения).
\end{itemize}

\textbf{Вывод:} Анализ показал, что существующие методы либо архаичны (живая очередь), либо не 
обеспечивают необходимого уровня безопасности и аналитики (таблицы/боты). Разработка специализированного 
веб-приложения на стеке Python/Flask позволит объединить преимущества облачных сервисов (доступность) с надежностью 
реляционных баз данных и специфической логикой учебного процесса (учет номеров лабораторных, статусов проверки и 
времени защиты). Это подтверждает целесообразность и актуальность выбранной темы курсовой работы.

\subsection{Обоснование выбора стека технологий}

Выбор программного стека (Solution Stack) является определяющим 
этапом проектирования, так как он влияет на масштабируемость системы, скорость разработки 
и требования к аппаратному обеспечению серверной части. Для реализации веб-приложения управления 
студенческими очередями был выбран стек на базе языка Python, ориентированный на высокую скорость 
обработки запросов и гибкость настройки.

\subsubsection{Выбор клиентских технологий}

Клиентская часть приложения (Frontend) отвечает за визуализацию данных и непосредственное взаимодействие 
пользователя с бизнес-логикой системы. Учитывая специфику эксплуатации приложения в условиях учебных корпусов 
(высокая нагрузка на базовые станции сотовой связи, нестабильный мобильный интернет и использование бюджетных 
мобильных устройств), при выборе стека основной приоритет был отдан минимизации веса страниц и максимальной 
скорости рендеринга.

\begin{itemize}
	\item \textbf{Сравнительный анализ и обоснование отказа от SPA-фреймворков:} 
В процессе проектирования рассматривалась возможность использования современных реактивных фреймворков, 
таких как \textbf{React, Vue.js и Angular}. Несмотря на их востребованность в индустрии, в рамках данного 
проекта от них было решено отказаться по ряду причин:

\begin{enumerate}
    \item \textit{Избыточность архитектуры:} Приложения на React или Angular требуют создания отдельного API 
	(REST или GraphQL). Для системы управления очередями, где количество сущностей ограничено, такая архитектура 
	удваивает объем работы без прямой выгоды для функционала.
    \item \textit{Производительность на мобильных устройствах:} SPA-фреймворки требуют времени на загрузку и 
	инициализацию JavaScript-бандла (Time to Interactive). В условиях нестабильного интернета в вузе задержка 
	в 5–10 секунд при загрузке тяжелого фреймворка неприемлема. Подход SSR (Flask + Jinja2) отправляет пользователю
	 уже готовый HTML-код, который отображается мгновенно.
    \item \textit{Сложность сборки и поддержки:} Использование данных технологий требует настройки инструментов 
	сборки (Vite, Webpack), управления Node.js-зависимостями и решения проблем с маршрутизацией на стороне клиента. В рамках учебного проекта это смещает фокус с разработки алгоритма очереди на решение инфраструктурных задач фронтенда.
\end{enumerate}

\item \textbf{HTML5 и CSS3:} Использование современных стандартов разметки позволяет создать семантически
 корректную структуру. Для реализации адаптивности под различные типы устройств (от смартфонов студентов до 
 широкоформатных мониторов преподавателей) применяются технологии Flexbox и Media Queries. Это критически важно, 
 так как интерфейс «Очереди» должен сохранять функциональность при любом разрешении экрана.

\item \textbf{Язык JavaScript (Native JS):} Для реализации динамических элементов интерфейса (например, 
обновления позиции в очереди в реальном времени или отправки формы подтверждения записи) выбран чистый 
JavaScript (ES6+). Использование Native JS вместо тяжелых библиотек позволяет избежать загрузки мегабайтов 
лишнего кода, что критично при работе через медленное соединение. Взаимодействие с сервером реализуется через 
современные стандарты Fetch API, что обеспечивает плавность работы без полной перезагрузки страницы.

\item \textbf{Шаблонизатор Jinja2:} Выбор данного инструмента обусловлен его бесшовной интеграцией с фреймворком 
Flask. Jinja2 реализует концепцию серверного рендеринга (SSR), что дает следующие преимущества:
\begin{itemize}
    \item создание иерархической структуры шаблонов через механизмы наследования (\texttt{block}, \texttt{extends}), 
	что упрощает поддержку кода;
    \item высокий уровень безопасности благодаря автоматическому экранированию данных (защита от XSS-инъекций);
    \item возможность генерации персонализированного контента (например, отображение кнопок управления только 
	для роли «Преподаватель») на стороне сервера, что исключает утечку логики управления в клиентский код.
\end{itemize}

Таким образом, комбинация \textbf{Jinja2 + Native JS} является наиболее рациональным выбором, обеспечивающим 
баланс между функциональностью и скоростью работы приложения в специфической учебной среде.
\end{itemize}

\subsubsection{Выбор серверных технологий и СУБД}

Серверная часть (Backend) является ядром системы, обеспечивающим бизнес-логику работы очереди и безопасность данных.

\begin{itemize}
 \item \textbf{Язык программирования Python:} Выбор Python обусловлен его статусом стандарта в области обработки 
 данных и веб-разработки. В сравнении с языками PHP или Java, Python обладает более лаконичным синтаксисом, что снижает 
 вероятность возникновения ошибок в логике приложения. Наличие огромного количества библиотек для работы с базами данных
  и аутентификацией делает его оптимальным для быстрой разработки (Rapid Application Development).
 
 \item \textbf{Микрофреймворк Flask:} Для реализации серверной логики был выбран Flask вместо его 
 основного конкурента — Django. Сравнительный анализ представлен в таблице~\ref{tab:flask_vs_django}.
\end{itemize}

\begin{table}[ht!]
 \centering
 \caption{Сравнение веб-фреймворков Flask и Django}
 \label{tab:flask_vs_django}
 \begin{tabular}{lp{3cm}p{3cm}}
  \toprule
  \textbf{Критерий} & \textbf{Flask} & \textbf{Django} \\
  \midrule
  Философия         & Микрофреймворк (ничего лишнего) & «Всё включено» (Monolith) \\
  Гибкость          & Максимальная (выбор компонентов за разработчиком) & Низкая (жёсткая структура) \\
  Порог вхождения   & Низкий & Средний/Высокий \\
  Производительность & Высокая за счёт минимализма & Зависит от количества подключенных модулей \\
  \bottomrule
 \end{tabular}
\end{table}

\textbf{Обоснование выбора Flask:} Поскольку проект предполагает создание специфической 
логики очереди с нуля, Flask предоставляет необходимую свободу в архитектурных решениях. Он не навязывает 
встроенную систему администрирования или специфические форматы форм, позволяя реализовать проект в строгом соответствии 
с требованиями технического задания.

\begin{itemize}
 \item \textbf{ORM SQLAlchemy:} Работа с базой данных осуществляется через Object-Relational Mapping 
 (ORM) библиотеку SQLAlchemy.
 \begin{itemize}
  \item \textit{Безопасность:} Автоматическая защита от SQL-инъекций.
  \item \textit{Удобство:} Работа с записями в БД как с объектами Python, что упрощает манипуляции 
  со списками студентов в очереди.
 \item \textbf{СУБД PostgreSQL:} В качестве системы управления базами данных выбрана реляционная СУБД PostgreSQL.
  \end{itemize}
  
 \begin{itemize}
  \item \textit{В сравнении с SQLite:} SQLite не поддерживает полноценную многопользовательскую запись 
  (Database Lock), что недопустимо в приложении, где десятки студентов могут одновременно пытаться встать в очередь.
  \item \textit{В сравнении с MySQL:} PostgreSQL обладает более строгим соответствием стандартам SQL 
  и лучшей поддержкой сложных транзакций, что гарантирует целостность очереди даже при одновременных запросах.
 \end{itemize}
\end{itemize}

\textbf{Вывод по подразделу 1.2:} Выбранный стек технологий (Python, Flask, PostgreSQL, Jinja2) представляет собой 
сбалансированное решение, позволяющее создать надежное, быстрое и легковесное веб-приложение. Использование микрофреймворка
 Flask в сочетании с мощной СУБД PostgreSQL обеспечивает необходимый фундамент для реализации сложной логики управления учебным 
 процессом в реальном времени.

 \subsubsection{Научная новизна и практическая значимость}

\textbf{Научная новизна} работы заключается в следующем:

\begin{itemize}
	\item Разработан авторский адаптивный алгоритм очереди, который, в отличие от стандартных FIFO-систем, 
	позволяет учитывать динамические параметры защиты: меняться местами с другими студнтами во время самого занятия,
	ставить отметки что работу студент сдал, отсутствие путаницы во время сдачи самих работ, так как все студенты
	видят количество сданнных работ друг-друга
	\item Предложена интегрированная модель учета метрик эффективности: система не просто фиксирует порядок, но и
	 собирает данные о среднем времени защиты одной работы, что в дальнейшем может быть использовано для 
	 оптимизации учебных планов и нагрузки преподавателей.
	\item Логика приложения создается «с нуля», что позволяет исключить избыточность универсальных 
	систем и достичь максимальной производительности интерфейса.
\end{itemize}

\textbf{Практическая значимость} работы состоит в создании готового к эксплуатации инструмента, 
который решает ряд острых проблем учебной группы:

\begin{itemize}
	\item \textit{Снижение конфликтности:} Исключение человеческого фактора и субъективности в определении 
	очередности.
	\item \textit{Прозрачность процесса:} Студент может в режиме реального времени видеть свою 
	позицию и прогнозировать время своего ответа, что позволяет эффективно распоряжаться временем вне аудитории.
	\item \textit{Информационная поддержка преподавателя:} Преподаватель может видеть сколько студентов на даннный
	момент присутствуют в аудитории, либо которые будут ему сдавать работы, из за этого он может учитывать то,
	сколько времени ему можно будет уделить студентам, чтобы все успеди сдать вовремя
\end{itemize}

\subsubsection{Описание архитектурных решений и требований к БД}

В рамках реализации проекта будет применена \textbf{трехуровневая архитектура}:

\begin{enumerate}
 \item \textit{Уровень представления (Presentation Layer):} Адаптивный веб-интерфейс, 
 обеспечивающий ввод данных и визуализацию текущего состояния очереди.
 \item \textit{Уровень бизнес-логики (Application Layer):} Обработка запросов на языке Python, 
 управление сессиями пользователей, реализация алгоритмов перестроения очереди.
 \item \textit{Уровень данных (Data Layer):} Реляционная СУБД PostgreSQL, обеспечивающая надежное 
 хранение информации о пользователях, истории защит и текущих активных сессиях.
\end{enumerate}

\textbf{Основные требования к базе данных} включают в себя обеспечение целостности связей 
(например, невозможность существования записи в очереди без привязки к конкретному студенту) и
 высокую скорость выполнения транзакционных запросов при одновременном обращении нескольких десятков пользователей.

